\documentclass[12pt,a4paper]{report}
\usepackage[]{amsmath,amsthm,amssymb,amscd}
\usepackage[latin1]{inputenc}
%\usepackage[spanish]{babel}
%\usepackage{fullpage}

\topmargin 0pt
\advance \topmargin by -\headheight
\advance \topmargin by -\headsep
     
\textheight 246mm
     
\oddsidemargin 0pt
\evensidemargin \oddsidemargin
\marginparwidth 0.5in
     
\textwidth 159mm


%               Theorem Environments
\theoremstyle{definition}
\newtheorem{ex}{}
\newcommand{\pd}[1]{\frac{\partial}{\partial #1}} %\pd{x}
%               Subquestions numbered with letters
\renewcommand{\theenumi}{\textbf{\alph{enumi}}}

%               Maths definitions

\newcommand{\NN}{\ensuremath{\mathbb N}}
\newcommand{\ZZ}{\ensuremath{\mathbb Z}}
\newcommand{\CC}{\ensuremath{\mathbb C}}
\newcommand{\RR}{\ensuremath{\mathbb R}}
\newcommand{\QQ}{\ensuremath{\mathbb Q}}
\newcommand{\g}{\ensuremath{\frak{g}}}
\newcommand{\h}{\ensuremath{\frak{h}}}
\newcommand{\ft}{\ensuremath{\frak{t}}}
\newcommand{\cB}{\mathcal{B}}
\newcommand{\cN}{\mathcal{N}}
\newcommand{\cL}{\mathcal{L}}
\newcommand{\cD}{\mathcal{D}}
%\newcommand{\pd}[1]{\frac{\partial}{\partial #1}} %\pd{x}
%\newcommand{\pd}[1]{\frac{\partial}{\partial #1}} %\pd{x}

\begin{document}
\noindent
{\sffamily\large Marco Zambon 
\hfill   a discutir:   21/02/2012\\Geometr\'ia III, curso 2011-12}
 

\noindent
\hrulefill{}\\
\begin{center}\bfseries\large\textsf{Hoja 1    \vspace{-2mm}}
\end{center}

\noindent
%\hrulefill{}\\


\begin{ex}  
Considera $$M:=\{(x,0):x <0\}\cup\{(x,y): x\ge0, y \in \RR\}\subset \RR^2$$
con la topolog\'ia subespacio inducida por $\RR^2$.
\begin{enumerate}
\item Demuestra que $M$ no es una variedad topol\'ogica.
\item Demuestra que $M$ no es una variedad topol\'ogica con borde.
\end{enumerate}

\textbf{Consejo:} Utiliza el teorema de invariancia del dominio.
\end{ex}

 \begin{ex}Sea $M$ una variedad topol\'ogica de dimension $n$ con borde, $p\in M$ y $(U,\phi)$ una carta con $p\in U$. 
 
Demuestra: si $\phi(p)\in R^{n-1}\times \{0\}$, entonces para toda otra carta $(U',\phi')$  con $p\in U'$, tenemos $\phi'(p)\in R^{n-1}\times \{0\}$.

 
\textbf{Consejo:} Utiliza el teorema de invariancia del dominio.
\end{ex}


 \begin{ex}Sea $M$ una variedad topol\'ogica de dimension $n$ con borde. Demuestra:
\begin{enumerate}
\item $Int(M)$ es una variedad topol\'ogica de dimension $n$ (sin borde).
\item $\partial M$ es una variedad topol\'ogica de dimension $n-1$ (sin borde).
\end{enumerate}
\end{ex}
 

\begin{ex}
\begin{enumerate}
\item Sea $X$ un espacio topol\'ogico compacto, $Y$ un espacio topol\'ogico Hausdorff, y $f \colon X \to Y$ una biyecci\'on continua. Demuestra: entonces $f$ es un homeomorfismo.
\item Sea $Z$ un espacio topol\'ogico y $\sim$ una relaci\'on de equivalencia en $Z$. Denota por $\pi \colon Z \to Z/\sim$ la aplicaci\'on canonica al espacio cociente. Describe las condiciones necesarias y suficientes para que una aplicaci\'on $F \colon Z \to Y$ induzca una aplicaci\'on  $\tilde{F} \colon Z/\sim \to Y$ (es decir, para que exista una aplicaci\'on $\tilde{F}$ tal que $F=\tilde{F}\circ\pi$). En este caso, demuestra que $F$ es continua si y solo si $\tilde{F}$ es continua. 
\item Sea $Z$ un espacio topol\'ogico compacto, $\sim$ una una relaci\'on de equivalencia en $Z$, y $Y$   un espacio topol\'ogico Hausdorff. Sea $F \colon Z \to X$ una aplicaci\'on continua, sobreyectiva, y tal que $z\sim z' \Leftrightarrow F(z)=F(z')$. Demuestra: $F$ induce un homeomorfismo $\tilde{F} \colon Z/\sim 
\to Y$. 
\end{enumerate}
\textbf{Nota:} Para a. es \'util saber que un subconjunto compacto de un espacio topologico Hausdorff es automaticamente cerrado.
\end{ex}





\begin{ex}
\begin{enumerate}
\item Sea $P_1(\RR)$ el cociente de $S^1:=\{z\in \CC: |z|=1\}$ por la relaci\'on de equivalencia $z\sim -z$. A qu\'e es homeomorfo $P_1(\RR)$?
\item Sea $M:= P_2(\RR)$, definida como en clase. El borde $\partial M$ est\'a vacio o no? Demuestra.
\item $P_2(\RR)$ es compacto o no? Demuestra.\end{enumerate}
\end{ex}





%\begin{ex}  
%Demuestra que 
%$$T \;\;\sharp \;\; C\cong T-(D_1 \cup D_2)$$
%donde $T=S^1\times S^1$ es el toro, $C=S^1\times [0,1]$ el cilindro, y $D_i\subset T$ (por $i=1,2$) son discos cerrados tal que $D_1\cap D_2=\emptyset$.
%\end{ex} 






 \end{document}
 %%%%%%%
 \begin{ex}

\end{ex}

 \begin{ex}
\begin{enumerate}
\item 
\end{enumerate}
\end{ex}
 
 
 
 